\documentclass[11pt]{article}
\usepackage[utf8]{inputenc}
\usepackage[T1]{fontenc}
\usepackage{amsmath,amssymb,mathtools}
\usepackage[margin=1in]{geometry}
\usepackage{hyperref}
\usepackage{array}
\usepackage{parskip}
\setlength{\parindent}{0pt}
\hypersetup{colorlinks=true,linkcolor=blue,urlcolor=blue}
\title{Lecture 14: Magnitude-Phase Representation and Filters}
\author{}
\date{}
\begin{document}
\maketitle

\section*{Polar Form of FT}

\[X(j\omega) = |X(j\omega)| \, e^{j\angle X(j\omega)}\]

\begin{itemize}
  \item \textbf{Magnitude spectrum} $|X(j\omega)|$: amplitude at each frequency
  \item \textbf{Phase spectrum} $\angle X(j\omega)$: phase shift at each frequency
\end{itemize}

\section*{Decibel Scale}

\[|H(j\omega)|_{\text{dB}} = 20 \log_{10} |H(j\omega)|\]

\begin{center}
\begin{tabular}{|l|l|}
\hline
\textbf{Linear} & \textbf{dB} \\
\hline
1 & 0 dB \\
\hline
$1/\sqrt{2} \approx 0.707$ & -3 dB \\
\hline
0.5 & -6 dB \\
\hline
0.1 & -20 dB \\
\hline
0.01 & -40 dB \\
\hline
10 & +20 dB \\
\hline
\end{tabular}
\end{center}

\textbf{Cutoff frequency:} defined at the \textbf{-3 dB point} where power is halved ($|H|^2 = \frac{1}{2}|H(0)|^2$).

\section*{Energy Spectral Density}

\[E = \int_{-\infty}^{\infty} |x(t)|^2 \, dt = \frac{1}{2\pi} \int_{-\infty}^{\infty} |X(j\omega)|^2 \, d\omega\]

$|X(j\omega)|^2$ tells how energy is distributed across frequency. Phase does not affect total energy.

\section*{Phase and Delay}

\subsection*{Linear Phase}

\[\angle H(j\omega) = -\omega t_0\]

A system with linear phase imposes a \textbf{pure time delay} of $t_0$ seconds on all frequencies equally. The output shape is preserved -- no distortion.

\[h(t) \text{ linear phase} \implies y(t) = c \cdot x(t - t_0)\]

\subsection*{Phase Distortion}

Nonlinear phase means different frequencies experience different delays. The output waveform shape is \textbf{distorted} even if magnitude is flat.

\subsection*{Group Delay}

\[\tau(\omega) = -\frac{d}{d\omega} \angle H(j\omega)\]

\begin{itemize}
  \item For linear phase: $\tau(\omega) = t_0$ (constant for all $\omega$)
  \item For nonlinear phase: $\tau(\omega)$ varies with frequency
\end{itemize}

\section*{Ideal Filters}

\subsection*{Ideal Lowpass Filter}

\[H(j\omega) = \begin{cases} 1, & |\omega| \leq \omega_c \\ 0, & |\omega| > \omega_c \end{cases}\]

Impulse response: $h(t) = \frac{\sin(\omega_c t)}{\pi t} = \frac{\omega_c}{\pi} \text{sinc}\left(\frac{\omega_c t}{\pi}\right)$

\textbf{Non-causal:} $h(t) \neq 0$ for $t < 0$, so ideal filters are \textbf{physically unrealizable}.

\subsection*{With Linear Phase (Delayed Ideal Filter)}

\[H(j\omega) = e^{-j\omega t_0}, \quad |\omega| \leq \omega_c\]

\[h(t) = \frac{\sin(\omega_c(t - t_0))}{\pi(t - t_0)}\]

Still noncausal for any finite $t_0$ (sinc extends to $-\infty$).

\section*{Nonideal (Real) Filter Specifications}

\begin{itemize}
  \item \textbf{Passband edge} $\omega_p$: end of the passband
  \item \textbf{Stopband edge} $\omega_s$: start of the stopband
  \item \textbf{Transition band}: $\omega_p < \omega < \omega_s$ (gradual roll-off region)
  \item \textbf{Passband ripple} $\delta_1$: maximum deviation from 1 in passband
  \item \textbf{Stopband attenuation} $\delta_2$: maximum leakage in stopband
\end{itemize}

\section*{Bandwidth and Rise Time Trade-off}

\[\omega_c \cdot t_r \approx \text{constant}\]

\begin{itemize}
  \item \textbf{Wider bandwidth} (larger $\omega_c$) $\Rightarrow$ \textbf{shorter rise time} (faster response)
  \item \textbf{Narrower bandwidth} $\Rightarrow$ \textbf{longer rise time} (slower, smoother response)
  \item Cannot have both sharp frequency cutoff and fast time response
\end{itemize}

\section*{Real-Part Sufficiency}

For a \textbf{real and causal} system, $H(j\omega)$ is completely determined by:
\begin{itemize}
  \item Its real part alone, OR
  \item Its imaginary part alone
\end{itemize}

The real and imaginary parts are related by the \textbf{Hilbert transform}.

\end{document}
